\documentclass[12pt,a4paper,oneside]{report}
\usepackage[T1]{fontenc}
\usepackage[utf8]{inputenc}
\usepackage[french]{babel}
\usepackage{lmodern,microtype,csquotes}
\usepackage[a4paper,top=2.5cm,bottom=2.5cm,left=2.8cm,right=2.5cm,headheight=15pt]{geometry}
\usepackage{xcolor,graphicx,float,booktabs,longtable,array,enumitem,fancyhdr,titlesec}
\graphicspath{{./}}
\usepackage{listings}
\usepackage[most]{tcolorbox}
\usepackage[hidelinks]{hyperref}

\definecolor{teal}{HTML}{0F766E}
\definecolor{tealclair}{HTML}{CCFBF1}
\definecolor{zinc}{HTML}{27272A}
\definecolor{gris}{HTML}{71717A}
\definecolor{papier}{HTML}{F8FAFC}
\definecolor{codebg}{HTML}{F4F4F5}
\definecolor{rouge}{HTML}{B91C1C}

\hypersetup{pdftitle={SecWatch — Rapport de projet},pdfauthor={Nom de l'étudiant},colorlinks=true,linkcolor=teal,urlcolor=teal}
\setlength{\parindent}{0pt}
\setlength{\parskip}{0.55em}
\setlist[itemize]{leftmargin=1.5em,itemsep=0.25em}
\setlist[enumerate]{leftmargin=1.7em,itemsep=0.3em}

\titleformat{\chapter}[display]
  {\normalfont\bfseries\color{zinc}}
  {\filleft\large\color{teal}\chaptername\ \thechapter}
  {0.8ex}
  {\titlerule\vspace{1.2ex}\Huge}
  [\vspace{0.8ex}\titlerule]
\titleformat{\section}{\Large\bfseries\color{zinc}}{\thesection}{0.8em}{}
\titleformat{\subsection}{\large\bfseries\color{teal}}{\thesubsection}{0.8em}{}

\pagestyle{fancy}
\fancyhf{}
\fancyhead[L]{\small\color{gris}\nouppercase{\leftmark}}
\fancyhead[R]{}
\fancyfoot[L]{\small\color{gris}SecWatch}
\fancyfoot[R]{\small\color{gris}\thepage}
\renewcommand{\headrulewidth}{0.4pt}
\renewcommand{\footrulewidth}{0pt}
\fancypagestyle{plain}{\fancyhf{}\fancyfoot[R]{\small\color{gris}\thepage}\renewcommand{\headrulewidth}{0pt}}

\lstdefinestyle{secwatch}{
  backgroundcolor=\color{codebg},basicstyle=\ttfamily\small,
  keywordstyle=\color{teal}\bfseries,commentstyle=\color{gris}\itshape,
  stringstyle=\color{rouge},numbers=left,numberstyle=\tiny\color{gris},
  numbersep=8pt,frame=single,rulecolor=\color{tealclair},framerule=0.6pt,
  breaklines=true,showstringspaces=false,tabsize=2,columns=fullflexible,
  keepspaces=true,upquote=true,xleftmargin=1.4em,framexleftmargin=1em
}
\lstset{style=secwatch}

% Les fichiers PNG référencés sont placés à la racine du projet Overleaf.
\newcommand{\capture}[2]{
  \begin{figure}[H]
    \centering
    \includegraphics[width=\textwidth,height=0.68\textheight,keepaspectratio]{#2.png}
    \caption{#1}
    \label{fig:#2}
  \end{figure}
}

\begin{document}

% PAGE DE GARDE
\begin{titlepage}
  \thispagestyle{empty}
  \centering
  \vspace*{1.3cm}
  {\large\color{teal}\textbf{RAPPORT DE PROJET}}\\[2.6cm]
  {\Huge\bfseries\color{zinc}SecWatch}\\[0.6cm]
  {\Large\color{gris}Application de déclaration et de suivi des incidents}\\[0.4cm]
  {\normalsize\color{gris}Sécurité des systèmes distribués}\\[1.8cm]
  \rule{0.72\textwidth}{0.8pt}\\[0.7cm]
  {\large Spring Boot \quad React \quad JWT}\\[0.7cm]
  \rule{0.72\textwidth}{0.8pt}
  \vfill
  \begin{tabular}{@{}ll@{}}
    \textbf{Étudiant·e :} & \rule{8cm}{0.4pt}\\[0.55cm]
    \textbf{Formation :} & \rule{8cm}{0.4pt}\\[0.55cm]
    \textbf{Enseignant·e :} & \rule{8cm}{0.4pt}\\[0.55cm]
    \textbf{Établissement :} & \rule{8cm}{0.4pt}\\[0.55cm]
    \textbf{Année universitaire :} & 2026--2027
  \end{tabular}
  \vspace*{1.5cm}
\end{titlepage}

\pagenumbering{roman}
\chapter*{Résumé}
SecWatch est une application web de déclaration et de suivi des incidents informatiques ou de cybersécurité. Elle distingue trois rôles : EMPLOYEE, TECHNICIAN et ADMIN. Son backend Spring Boot expose une API REST sécurisée par Spring Security et des jetons JWT. Son frontend React adapte les pages au rôle connecté et communique avec l'API au moyen d'Axios.

L'application permet de créer et suivre des incidents, de les affecter à un technicien, de modifier leur statut, d'ajouter des notes et de consulter l'historique. Elle comprend aussi les profils, l'administration des utilisateurs, les validations, les erreurs API et les tableaux de bord par rôle.

\textbf{Mots-clés :} JWT, Spring Security, React, JPA, authentification, contrôle d'accès, cybersécurité.
\tableofcontents
\clearpage
\pagenumbering{arabic}

\chapter{Présentation du projet}
\section{Contexte}
Les entreprises dépendent de leurs postes de travail, de leurs applications et de leurs réseaux. Une infection par logiciel malveillant, une tentative d'hameçonnage, une panne réseau, un problème d'accès ou une fuite potentielle de données doit pouvoir être signalé, attribué et suivi jusqu'à sa résolution.

SecWatch centralise ces signalements dans deux projets distincts : un backend Spring Boot et un frontend React. Les données sont persistées avec JPA. H2 est utilisé par défaut en développement; MySQL peut être configuré avec les variables d'environnement.

\section{Objectifs}
\begin{itemize}
  \item authentifier les utilisateurs avec un JWT à durée limitée;
  \item contrôler les accès selon les rôles et la propriété des incidents;
  \item valider les données et produire des erreurs HTTP cohérentes;
  \item créer, affecter et suivre les incidents;
  \item conserver l'historique et les notes de traitement;
  \item fournir des tableaux de bord adaptés aux rôles.
\end{itemize}

\section{Rôles}
\begin{longtable}{p{3cm}p{10.5cm}}
\toprule
\textbf{Rôle} & \textbf{Responsabilités}\\
\midrule
EMPLOYEE & Consulter et modifier son profil, créer des incidents et consulter uniquement ses propres signalements.\\
TECHNICIAN & Consulter les incidents affectés, modifier leur statut et ajouter des notes de traitement.\\
ADMIN & Consulter les incidents et utilisateurs, affecter les techniciens, activer ou désactiver les comptes et consulter les statistiques.\\
\bottomrule
\end{longtable}

\section{Périmètre}
L'inscription publique attribue toujours le rôle EMPLOYEE. Les comptes ADMIN et TECHNICIAN de démonstration sont préchargés. Le Refresh Token est facultatif dans le sujet et n'est pas implémenté; le frontend demande une nouvelle connexion lorsque l'access token expire.

\chapter{Architecture et organisation}
\section{Architecture générale}
Le navigateur charge React. Axios envoie les requêtes HTTP et ajoute le JWT à l'en-tête \texttt{Authorization}. Spring Security filtre les requêtes avant les contrôleurs. Les contrôleurs valident les entrées, les services appliquent les règles métier et les repositories JPA communiquent avec la base.

\begin{figure}[H]
\centering
\includegraphics[width=\textwidth,height=0.68\textheight,keepaspectratio]{diagram-architecture.png}
\caption{Architecture logique de SecWatch}
\label{fig:architecture}
\end{figure}

\section{Organisation du dépôt}
\begin{itemize}
  \item \texttt{backend/} : modèles JPA, contrôleurs, sécurité, services et repositories;
  \item \texttt{frontend/} : composants React, routes, pages et services Axios;
  \item \texttt{docs/} : rapport et diagrammes Mermaid/SVG;
  \item \texttt{postman/} : collection de requêtes de démonstration.
\end{itemize}

\section{Modèle de données}
Un incident possède un créateur obligatoire et peut être affecté à un technicien. L'historique et les commentaires sont des entités séparées : un incident peut recevoir plusieurs événements et notes. Les relations utilisent principalement \texttt{@ManyToOne}. Le chargement différé limite les chargements inutiles; les requêtes chargent les références nécessaires avant de construire les DTOs.

\begin{figure}[H]
\centering
\includegraphics[width=\textwidth,height=0.72\textheight,keepaspectratio]{diagram-entities.png}
\caption{Entités principales et relations de SecWatch}
\label{fig:entites}
\end{figure}
Les incidents référencent leur créateur et, facultativement, le technicien affecté. Chaque événement et chaque note est lié à l'incident ainsi qu'à l'utilisateur qui l'a enregistré.

Les diagrammes intégrés au rapport sont les fichiers PNG \texttt{diagram-entities.png} et \texttt{diagram-architecture.png}.

\chapter{Modèle de données et règles métier}
\section{Utilisateur}
L'entité \texttt{User} contient l'identifiant, le nom complet, l'email, le mot de passe encodé, le rôle, l'état actif, la date de création et la version du token. L'email est unique. \texttt{tokenVersion} invalide les JWT émis avant une désactivation ou un changement de mot de passe.

\section{Incident}
L'entité \texttt{Incident} contient le titre, la description, la catégorie, la priorité, le statut, les dates de création et de mise à jour, le créateur, le technicien facultatif et un champ \texttt{@Version} pour le verrouillage optimiste.

\begin{longtable}{p{3cm}p{10.5cm}}
\toprule
\textbf{Champ} & \textbf{Valeurs}\\
\midrule
Catégorie & NETWORK, MALWARE, PHISHING, ACCESS, SOFTWARE, HARDWARE, DATA\_LEAK, OTHER\\
Priorité & LOW, MEDIUM, HIGH, CRITICAL\\
Statut & OPEN, ASSIGNED, IN\_PROGRESS, RESOLVED, CLOSED\\
\bottomrule
\end{longtable}

\section{Historique et commentaires}
Chaque action importante est enregistrée dans \texttt{IncidentHistory} avec son acteur, son message et son horodatage. Les affectations et changements de statut créent un événement. \texttt{IncidentComment} conserve les notes de traitement. Le détail affiche les événements par ordre chronologique.

\section{Règles métier}
\begin{enumerate}
  \item Une inscription publique attribue EMPLOYEE; le client ne choisit jamais son rôle.
  \item Les mots de passe sont encodés avec BCrypt.
  \item Un employé ne consulte que ses propres incidents.
  \item Un technicien ne traite que les incidents qui lui sont affectés.
  \item Seul ADMIN affecte un technicien et consulte les données globales.
  \item Un compte désactivé ne peut plus se connecter.
  \item Un incident CLOSED ne change plus d'état et ne reçoit plus de note.
\end{enumerate}

\chapter{Authentification et sécurité}
\section{Inscription et connexion}
Les routes publiques \texttt{POST /auth/register} et \texttt{POST /auth/login} reçoivent des DTOs validés. Le DTO d'inscription ne contient pas de rôle. Le service encode le mot de passe avec BCrypt et affecte lui-même le rôle EMPLOYEE. La réponse de connexion ne contient aucun mot de passe.

\begin{lstlisting}[language=Java,caption={DTO d'inscription et rôle fixé côté serveur}]
public record Register(
    @NotBlank @Size(max = 100) String fullName,
    @NotBlank @Email @Size(max = 254) String email,
    @NotBlank @Size(min = 8, max = 72) String password
) {}

// Dans le service
user.password = passwordEncoder.encode(request.password());
user.role = Role.EMPLOYEE;
users.save(user);
\end{lstlisting}

\section{JWT}
Le JWT est signé en HMAC. Il contient l'email comme sujet, les dates d'émission et d'expiration, le rôle, l'identifiant utilisateur et la version du compte. Sa durée par défaut est de 300 secondes. Le parseur vérifie la signature et l'expiration avant de rendre les claims disponibles.

\begin{lstlisting}[language=Java,caption={Claims inclus dans le JWT}]
return Jwts.builder()
    .subject(user.email)
    .issuedAt(Date.from(now))
    .expiration(Date.from(now.plusSeconds(ttl)))
    .claim("role", user.role.name())
    .claim("uid", user.id)
    .claim("ver", user.tokenVersion)
    .signWith(key)
    .compact();
\end{lstlisting}

Un claim est une donnée portée par le token. \texttt{role} sert aux autorisations Spring, \texttt{uid} retrouve le compte même si l'email change et \texttt{ver} invalide les tokens antérieurs à une révocation.

\section{Filtre JWT et sessions sans état}
\texttt{JwtFilter} lit l'en-tête \texttt{Authorization: Bearer ...}, vérifie le token, recharge l'utilisateur, vérifie son activation et compare la version du JWT à celle enregistrée. Il place ensuite l'authentification dans le \texttt{SecurityContext}. \texttt{SessionCreationPolicy.STATELESS} indique que le serveur ne conserve pas de session HTTP.

\begin{lstlisting}[language=Java,caption={Extrait de la configuration Spring Security}]
http.csrf(csrf -> csrf.disable())
    .sessionManagement(session -> session
        .sessionCreationPolicy(SessionCreationPolicy.STATELESS))
    .authorizeHttpRequests(auth -> auth
        .requestMatchers("/auth/**", "/h2-console/**").permitAll()
        .requestMatchers("/api/**").authenticated()
        .anyRequest().permitAll());
\end{lstlisting}

Le CORS autorise le frontend, par défaut \texttt{http://localhost:5173}, ainsi que les méthodes et en-têtes nécessaires à Axios.

\section{Autorisations par rôle et propriétaire}
Les annotations \texttt{@PreAuthorize} contrôlent les rôles au niveau des endpoints. Le service vérifie également que l'utilisateur est le propriétaire de l'incident ou le technicien affecté. Le masquage d'un bouton React n'est pas une protection serveur.

\begin{lstlisting}[language=Java,caption={Contrôle d'accès sur un incident}]
boolean canSee(Incident incident, User user) {
    return user.role == Role.ADMIN
        || sameUser(incident.createdBy, user)
        || (user.role == Role.TECHNICIAN
            && sameUser(incident.assignedTo, user));
}
\end{lstlisting}

Les relations JPA différées sont lues par des getters ou des requêtes qui chargent les références nécessaires. Cela évite l'accès direct aux champs d'un proxy Hibernate non initialisé.

\section{Expiration et désactivation}
Quand une route protégée répond 401, Axios supprime le token et le profil local, puis redirige vers la connexion en affichant un message de session expirée. Un compte désactivé ne peut pas s'authentifier et ses anciens tokens sont rejetés. Le Refresh Token est facultatif et n'est pas inclus.

\chapter{API REST}
\section{Endpoints principaux}
\small
\begin{longtable}{p{2cm}p{7.2cm}p{4.3cm}}
\toprule
\textbf{Méthode} & \textbf{Endpoint} & \textbf{Accès}\\
\midrule
POST & \texttt{/auth/register}, \texttt{/auth/login} & Public\\
GET / PUT & \texttt{/api/profile} & Authentifié\\
PUT & \texttt{/api/profile/password} & Authentifié\\
POST & \texttt{/api/incidents} & EMPLOYEE\\
GET & \texttt{/api/incidents/my} & EMPLOYEE\\
GET & \texttt{/api/incidents/\{id\}} & Propriétaire, technicien affecté ou ADMIN\\
GET & \texttt{/api/incidents/assigned} & TECHNICIAN\\
GET & \texttt{/api/incidents} & ADMIN\\
PUT & \texttt{/api/incidents/\{id\}/status} & Technicien affecté ou ADMIN\\
PUT & \texttt{/api/incidents/\{id\}/assign/\{technicianId\}} & ADMIN\\
POST & \texttt{/api/incidents/\{id\}/comments} & Technicien affecté ou ADMIN\\
GET & \texttt{/api/admin/users} & ADMIN\\
PUT & \texttt{/api/admin/users/\{id\}/enable} ou \texttt{/disable} & ADMIN\\
GET & \texttt{/api/admin/statistics} & ADMIN\\
GET & \texttt{/api/admin/dashboard} & ADMIN; statistiques et incidents\\
\bottomrule
\end{longtable}
\normalsize

\section{Validation et codes HTTP}
Les DTOs utilisent \texttt{@NotBlank}, \texttt{@Email}, \texttt{@Size} et \texttt{@NotNull}; les contrôleurs déclenchent la validation avec \texttt{@Valid}. \texttt{@ControllerAdvice} transforme les exceptions en réponses JSON contenant le code HTTP et un message.

\begin{longtable}{p{2cm}p{11cm}}
\toprule
\textbf{Code} & \textbf{Signification dans l'application}\\
\midrule
400 & Données invalides, ancien mot de passe incorrect ou confirmation différente.\\
401 & JWT absent, invalide ou expiré; identifiants incorrects; compte désactivé.\\
403 & Rôle insuffisant ou tentative d'accès à l'incident d'un autre employé.\\
404 & Incident, utilisateur ou technicien introuvable.\\
409 & Adresse email déjà utilisée ou conflit métier.\\
\bottomrule
\end{longtable}

\chapter{Application React}
\section{Routes et organisation}
\texttt{ProtectedRoute} vérifie la présence d'un token et \texttt{RoleRoute} limite certaines pages aux rôles autorisés. Ces composants guident l'utilisateur; l'API effectue le contrôle d'accès réel. Les pages couvrent la connexion, l'inscription, les dashboards, le profil, les incidents, les affectations et l'administration.

\section{Axios et stockage du JWT}
Après connexion, le client conserve le token dans le navigateur. L'intercepteur Axios l'ajoute aux requêtes protégées. Le gestionnaire de réponse traite l'expiration :

\begin{lstlisting}[language=JavaScript,caption={Intercepteurs Axios}]
api.interceptors.request.use((config) => {
  const token = localStorage.getItem("token");
  if (token) config.headers.Authorization = "Bearer " + token;
  return config;
});

api.interceptors.response.use((response) => response, (error) => {
  if (error.response?.status === 401 && localStorage.getItem("token")) {
    localStorage.removeItem("token");
    localStorage.removeItem("user");
    sessionStorage.setItem("session-expired", "1");
    window.location.href = "/login";
  }
  return Promise.reject(error);
});
\end{lstlisting}

\section{Tableaux de bord}
\begin{description}[leftmargin=2.8cm,style=nextline]
  \item[EMPLOYEE] Nombre d'incidents personnels, ouverts, en cours, résolus ou clôturés, et signalements récents.
  \item[TECHNICIAN] Incidents affectés, priorités critiques, incidents en cours et résolus.
  \item[ADMIN] Utilisateurs, employés, techniciens, incidents par statut, incidents critiques et tableau des incidents récents.
\end{description}

Le dashboard administrateur reçoit ses compteurs et ses lignes depuis un endpoint unique. Le frontend vérifie que le nombre reçu correspond au total annoncé. L'interface utilise une palette zinc et teal, un thème clair/sombre mémorisé localement et des composants inspirés de shadcn/ui. Les fichiers frontend sont formatés avec Prettier.

\chapter{Installation et démonstration}
\section{Prérequis}
Java 21, Maven 3.9 ou supérieur, Node.js et npm.

\section{Démarrage du backend}
\begin{lstlisting}[language=bash,caption={Lancer Spring Boot}]
cd backend
mvn spring-boot:run
\end{lstlisting}
L'API écoute sur \texttt{http://localhost:8080}. La console H2 est disponible à \texttt{http://localhost:8080/h2-console}. URL JDBC : \texttt{jdbc:h2:file:./.data/secwatch}; utilisateur \texttt{sa}, mot de passe vide.

\section{Démarrage du frontend}
\begin{lstlisting}[language=bash,caption={Lancer React}]
cd frontend
npm install
npm run dev
\end{lstlisting}
Vite affiche l'adresse locale, normalement \texttt{http://localhost:5173}. Si le port est occupé, ouvrir l'adresse réellement affichée.

\section{Comptes de démonstration}
Les comptes sont préchargés si \texttt{SEED\_DEMO=true}. Mot de passe local commun : \texttt{Demo1234!}.
\begin{center}
\begin{tabular}{ll}
\toprule
Rôle & Email\\
\midrule
ADMIN & \texttt{admin@secwatch.local}\\
TECHNICIAN & \texttt{tech@secwatch.local}\\
EMPLOYEE & \texttt{employee@secwatch.local}\\
EMPLOYEE & \texttt{employee2@secwatch.local}\\
\bottomrule
\end{tabular}
\end{center}

\section{Scénario fonctionnel}
\begin{enumerate}
  \item Se connecter comme premier employé et créer deux incidents.
  \item Se connecter comme second employé et créer un incident.
  \item Vérifier que le premier employé ne peut pas lire l'incident du second.
  \item En ADMIN, consulter les incidents et affecter l'un d'eux au technicien.
  \item En TECHNICIAN, passer l'incident à IN\_PROGRESS, ajouter une note puis le passer à RESOLVED.
  \item Revenir au compte employé et vérifier le statut et l'historique.
\end{enumerate}

\chapter{Vérifications de sécurité avec curl}
Les neuf situations ci-dessous couvrent les principaux contrôles d'accès et de validation de l'API.

\section{Cas requis}
\begin{longtable}{p{0.7cm}p{3.9cm}p{4.8cm}p{3.3cm}}
\toprule
\textbf{N°} & \textbf{Test} & \textbf{Requête / préparation} & \textbf{Attendu}\\
\midrule
1 & API sans JWT & GET \texttt{/api/profile}, sans Authorization & 401\\
2 & JWT invalide & GET profil avec \texttt{Bearer invalide} & 401\\
3 & JWT expiré & Rejouer le token après expiration & 401\\
4 & Employé sur route ADMIN & GET \texttt{/api/admin/users} avec token employé & 403\\
5 & Incident d'un autre employé & GET \texttt{/api/incidents/ID} avec mauvais propriétaire & 403\\
6 & Technicien non affecté & PUT statut d'un incident non affecté & 403\\
7 & Compte désactivé & POST login après désactivation & 401\\
8 & Email déjà utilisé & Répéter l'inscription avec le même email & 409\\
9 & Incident invalide & POST avec titre vide ou priorité absente & 400\\
\bottomrule
\end{longtable}

\section{Résultats des vérifications}
Les essais couvrent l'absence et l'invalidité du JWT, l'expiration, les rôles, la propriété des incidents, l'affectation, la désactivation, les doublons et la validation des données. Les réponses observées sont illustrées ci-dessous.

\capture{API sans JWT — réponse 401}{curl-sans-jwt}
\capture{JWT invalide — réponse 401}{curl-jwt-invalide}
\capture{JWT expiré — réponse 401}{curl-jwt-expire}
\capture{Accès employé à une route administrateur — réponse 403}{curl-route-admin}
\capture{Accès à l'incident d'un autre employé — réponse 403}{curl-proprietaire}
\capture{Technicien non affecté — réponse 403}{curl-technicien}
\capture{Connexion d'un compte désactivé — réponse 401}{curl-compte-desactive}
\capture{Inscription avec email déjà utilisé — réponse 409}{curl-email-duplique}
\capture{Incident invalide — réponse 400}{curl-validation}

\chapter{Captures de l'application}
Les figures suivantes présentent les écrans principaux de SecWatch.

\capture{Connexion}{ui-login}
\capture{Inscription}{ui-register}
\capture{Dashboard employé}{ui-dashboard-employee}
\capture{Détail d'un incident côté employé}{ui-incident-employee}
\capture{Dashboard administrateur}{ui-dashboard-admin}
\capture{Gestion des utilisateurs}{ui-admin-users}
\capture{Incidents affectés au technicien}{ui-assigned-incidents}
\capture{Détail et traitement d'un incident}{ui-technician-detail}

\chapter{Résultats, limites et conclusion}
\section{Résultats}
SecWatch met en œuvre l'authentification JWT, les rôles, la protection des routes, le contrôle du propriétaire, la validation, les erreurs HTTP, les profils et le suivi des incidents. L'affectation, les statuts, les notes et l'historique couvrent le parcours de traitement. Les tableaux de bord présentent des informations adaptées à chaque rôle.

\section{Limites connues}
Le Refresh Token, facultatif dans le sujet, n'est pas inclus. Les comptes de démonstration sont réservés au développement local.

\section{Conclusion}
Le projet met en pratique les mécanismes de sécurité d'une application distribuée. Le frontend facilite l'utilisation, mais le backend vérifie le JWT, le rôle et les droits sur chaque incident. La séparation entre contrôleurs, services, repositories et pages React clarifie les responsabilités du système.

\end{document}
